\chapter{Iron Supplementation}

\section*{Definition and Overview}
Iron supplementation is the use of iron medicines to treat or prevent iron deficiency. It is most commonly given as oral iron tablets, which are effective and inexpensive, and as intravenous iron when tablets are not tolerated, not absorbed, or needed quickly. Iron supplements are used for people with confirmed iron deficiency, for pregnant women to prevent deficiency, and for people with ongoing losses such as heavy periods. They should not be taken without a clear reason, because excess iron can be harmful. When prescribed appropriately, iron supplements restore energy, correct anaemia, and prevent the complications of deficiency.

\section*{Epidemiology}
Iron deficiency is the most common nutritional deficiency worldwide, and iron supplementation is one of the most widely used treatments in medicine. In Australia, around one in eight women of reproductive age is iron deficient, and iron is recommended routinely in pregnancy, where needs outstrip intake for most women. Supplements are also used for infants at risk, for people with gut conditions that impair absorption, and after surgery or blood loss. The use of intravenous iron has grown as newer, safer preparations have become available.

\section*{Aetiology and Pathophysiology}
Iron deficiency develops when intake does not match needs, from menstruation, pregnancy, growth, poor diet, malabsorption, or bleeding. Oral iron supplements deliver elemental iron that is absorbed in the upper small intestine, mainly as non-haem iron, and transported to the bone marrow for red blood cell production and to stores. Absorption is enhanced by vitamin C and reduced by tea, coffee, and calcium. Some people cannot tolerate oral iron because of nausea, constipation, or a metallic taste, and some absorb it poorly, such as those with inflammatory bowel disease or after gastric surgery. Intravenous iron bypasses the gut, delivering iron directly to the bloodstream in a single or several doses.

\section*{Clinical Features}
\begin{itemize}
    \item Oral iron is usually taken daily or every second day, ideally with vitamin C
    \item Side effects include nausea, constipation, dark stools, and a metallic taste
    \item Taking iron with food reduces side effects but also reduces absorption
    \item Iron is best absorbed away from tea, coffee, and dairy
    \item Intravenous iron is given in a clinic or hospital over minutes to an hour
    \item Symptoms of deficiency improve within weeks, with blood counts normalising over one to two months
    \item Iron is dangerous for children in overdose, so it must be stored out of reach
\end{itemize}

\section*{Differential Diagnosis}
The decision to supplement iron follows confirmation that iron deficiency is present.
\begin{itemize}
    \item \textbf{Confirmed iron deficiency:} low ferritin and anaemia on blood tests
    \item \textbf{Prevention in pregnancy:} routine supplementation where recommended
    \item \textbf{Anaemia of chronic disease:} iron is not the primary treatment and may be harmful if not deficient
    \item \textbf{Other deficiencies:} B12 or folate deficiency needs different treatment
    \item \textbf{Thalassaemia:} anaemia without iron deficiency should not be treated with iron
    \item \textbf{Iron overload:} supplementation is dangerous in haemochromatosis
\end{itemize}

\section*{When to Seek Medical Care}
Iron supplements should be taken under medical advice after a confirmed diagnosis, and people taking them should be reviewed to confirm response and adjust the dose. Anyone with severe side effects, worsening symptoms, or new bleeding should seek review. Pregnant women should follow their antenatal team's advice about iron. Iron overdose is a medical emergency, especially in children.

\textbf{Call 000 (or local emergency services) for:}
\begin{itemize}
    \item Suspected iron overdose, especially in a child: vomiting, abdominal pain, diarrhoea, or drowsiness
    \item Severe breathlessness, chest pain, or collapse
    \item Black, tarry, or bloody stools, or vomiting blood
    \item Heavy bleeding with faintness or dizziness
    \item Signs of a severe reaction to an intravenous iron infusion
\end{itemize}

\section*{Diagnosis and Investigations}
\begin{itemize}
    \item \textbf{Full blood count and iron studies:} confirming deficiency before treatment
    \item \textbf{Ferritin:} the baseline measure of iron stores
    \item \textbf{Assessment of the cause:} bleeding, diet, absorption, and pregnancy
    \item \textbf{Dose selection:} based on age, weight, and the degree of deficiency
    \item \textbf{Monitoring:} repeat blood tests at four to eight weeks
    \item \textbf{Review of tolerability:} adjusting the preparation or route as needed
    \item \textbf{Pregnancy care:} routine screening and supplementation
\end{itemize}

\section*{Management}
\begin{itemize}
    \item \textbf{Oral iron:} tablets or liquid, taken as directed, with dose adjustments to reduce side effects
    \item \textbf{Alternate-day dosing:} can improve absorption and reduce side effects
    \item \textbf{Vitamin C:} taken with iron to enhance absorption
    \item \textbf{Intravenous iron:} for intolerance, malabsorption, severe deficiency, or late pregnancy
    \item \textbf{Treatment of the cause:} heavy periods, bleeding, and gut disease
    \item \textbf{Dietary support:} iron-rich foods alongside supplementation
    \item \textbf{Follow-up:} confirming correction of haemoglobin and stores
    \item \textbf{Continuing care:} maintenance or monitoring as appropriate
\end{itemize}

\section*{Complications}
\begin{itemize}
    \item Gastrointestinal side effects, including constipation and nausea
    \item Iron overload from unnecessary or prolonged supplementation
    \item Interactions with other medicines, including some thyroid and stomach medicines
    \item Acute iron poisoning from overdose, especially in children
    \item Infusion reactions with intravenous iron
    \item Darkening of the teeth with liquid iron preparations
\end{itemize}

\section*{Prognosis and Outcomes}
Iron supplementation is highly effective when used correctly. Most people with iron deficiency feel better within weeks, and anaemia resolves within one to two months, with stores replenished over several further months. Intravenous iron corrects deficiency reliably even when oral iron fails. With treatment of the underlying cause, recurrence is uncommon. The main risks arise from inappropriate use, so supplementation should always be based on a confirmed need and monitored.

\section*{Prevention and Self-Care}
\begin{itemize}
    \item Take iron only under medical advice and at the prescribed dose
    \item Take iron with vitamin C and away from tea, coffee, and dairy
    \item Space doses to improve tolerance and absorption
    \item Store iron safely, out of the reach of children
    \item Report side effects rather than stopping treatment abruptly
    \item Attend follow-up blood tests
    \item Continue dietary iron intake during and after supplementation
    \item Follow antenatal advice on iron in pregnancy
    \item Seek immediate help for any suspected overdose
\end{itemize}

\section*{Sources and Further Reading}
The following reputable sources provide further information for healthcare students and patients seeking reliable, current advice about iron supplementation.
\begin{itemize}
    \item Healthdirect Australia. \textit{Iron supplements.} https://www.healthdirect.gov.au/
    \item The Royal Australian College of General Practitioners. \textit{Iron supplementation guidance.} https://www.racgp.org.au/
    \item National Health and Medical Research Council. \textit{Iron in pregnancy guidelines.} https://www.nhmrc.gov.au/
    \item Pregnancy, Birth and Baby. \textit{Iron and pregnancy.} https://www.pregnancybirthbaby.org.au/
    \item Poisons Information Centre. \textit{Iron poisoning information.} https://www.poisonsinfo.nsw.gov.au/
    \item NHS. \textit{Iron supplements.} https://www.nhs.uk/
\end{itemize}
